\exercisesection{}

\begin{enumerate}
\def\labelenumi{\arabic{enumi})}
\tightlist
\item
  设 \(\mathcal{T} = (\mathrm{K}_i, p_{ij})\) 是以集合 I 为指标的紧空间的一个逆系统，设 K 是一个紧空间，并设 \((p_i)_{i \in I}\) 是连续映射 \(p_i: \mathrm{K} \to \mathrm{K}_i\) 组成的一个相容且分离的族。
\end{enumerate}

\begin{enumerate}
\def\labelenumi{\alph{enumi})}
\item
  设 A 是 K 上形如 \(f_{i} \circ p_{i}\) 的连续函数所成的集合，其中 i 取遍 I，\(f_{i}\) 取遍 \(K_{i}\) 上连续实函数所成的集合。证明 A 是 \(\mathcal{C}(K)\) 的一个稠密线性子空间。
\item
  设 \((\mu_i)_{i \in I}\) 是 \(\mathcal{T}\) 上正测度组成的一个逆子系统。假设这些 \(p_i\) 都是满射。对每个 \(f \in A\)，用 \(I_f\) 表示这样的 \(i \in I\) 组成的集合：\(f\) 形如 \(f_i \circ p_i\)，其中 \(f_i \in \mathcal{C}(K_i)\)（必然唯一）。证明存在唯一的测度 \(\pi\)（在 \(K\) 上），使得 \(\pi(f) = \inf_{i \in I_f} \mu_i(f_i)\)（对所有 \(f \in A\)）。证明 \(\pi\) 是
\end{enumerate}

K 上正测度 \(\mu\)（满足 \(p_i(\mu) \leqslant \mu_i\)，对所有 \(i \in I\)）中最大的一个。

¶ 2) 假设与记号同 No.~2 的 Th. 1，我们打算对它给出一个新的证明。

\begin{enumerate}
\def\labelenumi{\alph{enumi})}
\item
  设 K 是 T 的一个紧子集；对每个 \(i \in I\)，置 \(\mathrm{K}_{i} = p_{i}(\mathrm{T})\)，并用 \(q_{i}\) 表示 K 到 \(K_{i}\) 上的、在 K 上与 \(p_{i}\) 重合的映射；对 \(i \leqslant j\)，用 \(q_{ij}\) 表示 \(K_{j}\) 到 \(K_{i}\) 中的、在 \(K_{j}\) 上与 \(p_{ij}\) 重合的映射。由上一习题推出，存在 K 上最大的正测度 \(\pi^{K}\)，使得 \(q_{i}(\pi^{K}) = (\mu_{i})_{\mathbf{K}_{i}}\)（对所有 \(i \in I\)）。
\item
  若 K 与 L 是 T 的紧子集且 \(K \subset L\)，证明 \((\pi^{\mathrm{L}})_{\mathrm{K}} \geqslant \pi^{\mathrm{K}}\)；由此推出，形如 \(i^{\mathrm{K}}(\pi^{\mathrm{K}})\) 的测度所成的集合（其中 K 取遍 T 的紧子集所成的集合，而 \(i^{K}\) 是 K 到 T 中的典范单射）有一个上确界 \(\mu\)（在 \(\mathcal{M}_{+}(\mathrm{T})\) 中）。
\item
  证明 \(p_i(\mu) = \mu_i\) 对每个 \(i \in I\) 成立（只须注意 \(p_i(\mu) \leqslant \mu_i\)，且由假设 (P)，测度 \(p_i(\mu)\) 与 \(\mu_i\) 有相同的总质量）。
\end{enumerate}

